% kinematic_channel.tex
% The kinematic directional-dispersion channel: closed-form axial/transverse
% group-speed contrast along the optical axis (1,1,-1).
% NOTE (terminology, external review 2026-07-17): "birefringence" is reserved
% for the gauge sector; this channel is directional dispersion anisotropy.
%
% STATUS: closed-form split DERIVED and numerically verified
% (notes/closed_form_speed_split.md; exp_02). One framing decision remains
% before flipping the audit row PART->PASS: which observable the
% "speed split" names (group-velocity vs dispersion-curvature) -- see the
% photon-limit remark below.  Cross-checks: exp_01 (observability), exp_02
% (analytic identities).

\label{sec:kinematic_channel}

A terminological caution frames this section. The effect below is a uniaxial
\emph{directional dispersion anisotropy}: it contrasts propagation along
$\hat{\mathbf{n}}_*$ with propagation across it. It is \emph{not} birefringence in
the optical sense---two polarization eigenmodes with different speeds at one fixed
wavevector---and we reserve that word for the gauge sector
(Sec.~\ref{sec:gauge_sector}). The two branches $\lambda_\pm(\mathbf{k})$ below are
the eigenvalue branches of the spinor transfer operator (the paired
positive/negative-frequency spinor modes), not two optical polarizations. A second
caution concerns interpretation: the transfer operator is \emph{non-unitary} away
from the axis ($|\lambda_\pm|^2<1$), so identifying the eigenphase gradient with a
physical group velocity needs an argument. The A=1 normalization supplies it for the
\emph{massive} sector but---as we show---\emph{not} in the massless limit, where the
eigenphase does not disperse at all (Sec.~\ref{subsec:a1_normalization},
Prop.~\ref{prop:massless_flat}). What is rigorous throughout is the exact
axial--transverse anisotropy of the structure-factor magnitude; the group-velocity
reading is earned for $\omega\neq0$, and the massless and photon readings are treated
with the care they require.

The kinematic channel inherits Paper~I's single-tick spinor
propagator~\cite{menendez2026geometryfirst}
\begin{equation}
    T(\mathbf{k}) = \begin{pmatrix}
        i\sin(\omega/2) & \cos(\omega/2)\,\tilde{H}_\text{RGB}(\mathbf{k}) \\
        \cos(\omega/2)\,\tilde{H}_\text{CMY}(\mathbf{k}) & i\sin(\omega/2)
    \end{pmatrix},
    \qquad
    \tilde{H}_\text{CMY} = \overline{\tilde{H}_\text{RGB}},
    \label{eq:propagator}
\end{equation}
whose eigenvalues
$\lambda_\pm(\mathbf{k}) = i\sin(\omega/2) \pm
\cos(\omega/2)\,|\tilde{H}_\text{RGB}(\mathbf{k})|$ carry the entire
$\mathbf{k}$-dependence through the single scalar
$|\tilde{H}_\text{RGB}|$.  Here $\omega$ is the per-tick Zitterbewegung
rate ($m = \sin(\omega/2)/a$), and the massless \emph{spinor} limit is
$\omega\to0$ (its identification with a photon is examined, and not asserted, in
Remark~\ref{rem:photon_limit}).  The RGB basis vectors are
$\mathbf{V}_1=(1,1,1)$, $\mathbf{V}_2=(1,-1,-1)$, $\mathbf{V}_3=(-1,1,-1)$,
with $\mathbf{V}_1+\mathbf{V}_2+\mathbf{V}_3=(1,1,-1)$ fixing the optical
axis $\hat{\mathbf{n}}_* = (1,1,-1)/\sqrt{3}$.

\subsection{Exact structure factor and the perpendicular projector}
\label{subsec:structure_factor}

With $\tilde{H}_\text{RGB}(\mathbf{k}) = \tfrac13\sum_{\mathbf{V}\in\text{RGB}}
e^{i\mathbf{k}\cdot\mathbf{V}}$ and the difference vectors
$\mathbf{V}_1-\mathbf{V}_2=(0,2,2)$, $\mathbf{V}_1-\mathbf{V}_3=(2,0,2)$,
$\mathbf{V}_2-\mathbf{V}_3=(2,-2,0)$, the modulus is closed-form to all
orders:
\begin{equation}
    |\tilde{H}_\text{RGB}(\mathbf{k})|^2
    = \tfrac19\Bigl[\,3
        + 2\cos 2(k_y{+}k_z)
        + 2\cos 2(k_x{+}k_z)
        + 2\cos 2(k_x{-}k_y)\,\Bigr].
    \label{eq:H2_exact}
\end{equation}
Along the optical axis $\mathbf{k}=t\,(1,1,-1)$ every argument vanishes, so
$|\tilde{H}_\text{RGB}|^2 = 1$ \emph{exactly for all $t$}: the dispersion
is rigorously flat along $\hat{\mathbf{n}}_*$.  Expanding to second order
recovers Paper~I's $M_\text{eff}$, which takes the transparent form
\begin{equation}
    |\tilde{H}_\text{RGB}(\mathbf{k})|^2 = 1 - \mathbf{k}^T M_\text{eff}\,\mathbf{k}
        + O(k^4),
    \qquad
    M_\text{eff} = \tfrac43\bigl(\mathbb{I} - \hat{\mathbf{n}}_*\hat{\mathbf{n}}_*^{T}\bigr)
        = \tfrac43 P_\perp .
    \label{eq:Meff_projector}
\end{equation}
That is, $M_\text{eff}$ is $\tfrac43$ times the projector onto the plane
perpendicular to the optical axis (eigenvalues $\{4/3,4/3,0\}$).  The
dispersion therefore depends \emph{only} on the perpendicular wavevector
$\mathbf{k}_\perp = P_\perp\mathbf{k}$:
\begin{equation}
    |\tilde{H}_\text{RGB}(\mathbf{k})|^2 = 1 - \tfrac43|\mathbf{k}_\perp|^2 + O(k^4).
    \label{eq:H2_perp}
\end{equation}

\subsection{The eigenphase dispersion and the speed split}
\label{subsec:speed_split}

The mode energy (phase advance per tick) is
$\theta(\mathbf{k}) = \arg\lambda_+(\mathbf{k})$.  From
Eq.~\eqref{eq:propagator}, $\mathrm{Re}\,\lambda_+
= \cos(\omega/2)|\tilde{H}_\text{RGB}|$ and
$\mathrm{Im}\,\lambda_+ = \sin(\omega/2)$, giving the exact dispersion
\begin{equation}
    \tan\theta(\mathbf{k}) = \frac{\tan(\omega/2)}{|\tilde{H}_\text{RGB}(\mathbf{k})|}.
    \label{eq:dispersion_exact}
\end{equation}
At $\mathbf{k}=0$ this returns $\theta=\omega/2$ (the rest-mass phase);
off-axis $|\tilde{H}_\text{RGB}|<1$ raises $\theta$.  Differentiating
gives the group velocity in closed form,
\begin{equation}
    \mathbf{v}_g(\mathbf{k}) = \nabla_{\mathbf{k}}\theta
    = -\,\frac{\tfrac12\sin\omega\;\nabla|\tilde{H}_\text{RGB}|}
             {\cos^2(\omega/2)\,|\tilde{H}_\text{RGB}|^2 + \sin^2(\omega/2)},
    \label{eq:group_velocity}
\end{equation}
which, using Eq.~\eqref{eq:H2_perp}, reduces at small $\mathbf{k}$ to
\begin{equation}
    \mathbf{v}_g(\mathbf{k}) = \tfrac23\sin\omega\;\mathbf{k}_\perp + O(k^3).
    \label{eq:vg_small_k}
\end{equation}

\begin{proposition}[Axial--transverse group-velocity contrast]
\label{prop:speed_split}
For a wavevector at angle $\alpha$ to the optical axis
($|\mathbf{k}_\perp| = |\mathbf{k}|\sin\alpha$), the kinematic group speed
(read from the eigenphase, subject to the interpretation caveat below) is, to
leading order,
\begin{equation}
    |\mathbf{v}_g| = \tfrac23\sin\omega\;|\mathbf{k}|\,\sin\alpha .
    \label{eq:speed_law}
\end{equation}
The transverse direction ($\mathbf{k}\perp\hat{\mathbf{n}}_*$, $\alpha=90^\circ$)
carries the maximal group speed $\tfrac23\sin\omega\,|\mathbf{k}|$; the axial
direction ($\mathbf{k}\parallel\hat{\mathbf{n}}_*$, $\alpha=0$) is frozen,
$|\mathbf{v}_g|=0$.  The axial--transverse contrast is
\begin{equation}
    \Delta v \equiv v_\perp - v_\parallel = \tfrac23\sin\omega\;|\mathbf{k}|,
    \label{eq:split}
\end{equation}
with $\hat{\mathbf{n}}_* = (1,1,-1)/\sqrt{3}$ the frozen (axial) direction.
This is a contrast between two propagation \emph{directions}, not a
polarization split at fixed $\mathbf{k}$.
\end{proposition}

\subsection{The A=1 normalization and the massless limit}
\label{subsec:a1_normalization}

Away from the axis the transfer operator is not unitary:
$|\lambda_\pm(\mathbf{k})|^2 = \sin^2(\omega/2) + \cos^2(\omega/2)\,
|\tilde H_\text{RGB}(\mathbf{k})|^2 < 1$, so each eigenvalue carries amplitude
change as well as phase, and $\nabla_{\mathbf{k}}\theta$ is not \emph{a priori} a
physical group velocity. The A=1 rule bears on this, with a caveat we state carefully.
The scheduler renormalizes the amplitude to unit norm every tick---a real, positive
rescaling $\psi\to\psi/\|\psi\|$ that leaves every phase untouched and writes back
$\phi=\arg\psi$. For a pure Bloch \emph{eigenmode} this removes the scalar attenuation
while preserving the eigenphase, replacing $\lambda_\pm$ by
$\lambda_\pm/|\lambda_\pm|$, so $\theta_\pm(\mathbf{k}) = \arg\lambda_\pm(\mathbf{k})$
is the physical single-mode dispersion. For a general state the renormalization is
\emph{nonlinear}---a single global norm does not equalize the different $|\lambda(\mathbf{k})|$
of a packet's several $\mathbf{k}$-components---so it is \emph{not} a unitary operator on
the Hilbert space, and the group-velocity reading must be justified at the level of the
wave packet rather than by calling the map unitary. For a narrow-band packet a
stationary-phase estimate gives the centroid velocity as the phase-gradient
$\nabla_{\mathbf{k}}\theta$ (Eqs.~\eqref{eq:dispersion_exact},
\eqref{eq:group_velocity}), while the momentum dependence of $|\lambda|$ contributes an
additional amplitude filtering/spreading term, treated separately below and at $\omega=0$.
This transport reading is corroborated by \texttt{exp\_01} under the full nonlinear
dynamics; the distinction between the massive-sector phase dispersion and the massless
magnitude filtering is exactly the one drawn next.

\begin{proposition}[Massless-limit degeneracy]
\label{prop:massless_flat}
At $\omega=0$ the eigenvalues $\lambda_\pm = \pm|\tilde H_\text{RGB}(\mathbf{k})|$ are
real, so $\theta_\pm(\mathbf{k}) \in \{0,\pi\}$ independently of $\mathbf{k}$: the
eigenphase is exactly flat and the group velocity vanishes in every direction. (This
is immediate from the closed form---$\lambda_\pm$ real makes $\arg\lambda_\pm$
constant---and a direct finite-difference sweep confirms
$\max_{\mathbf{k}}|\nabla_{\mathbf{k}}\theta_+| = 0$ at $\omega=0$, versus $O(1)$ for
$\omega\neq0$.)
\end{proposition}

The A=1 normalization does not restore a dispersion here---the phase is already
$\mathbf{k}$-independent---so at $\omega=0$ the entire anisotropy is carried by the
eigenvalue \emph{magnitude} $|\tilde H_\text{RGB}|$, and there is no eigenphase energy
dispersion. It is important not to over-read the single-mode step: fixing
$|\lambda_{\mathbf{k}}|\to1$ per plane wave would flatten one mode's magnitude, but the
physical A=1 rule renormalizes the \emph{total} norm of a state, which does \emph{not}
equalize the relative magnitudes across the $\mathbf{k}$-components of a wavepacket. The
magnitude anisotropy therefore survives into observable amplitudes---this is exactly the
$\omega$-independent anisotropic spreading that \texttt{exp\_01} measures (present and
identical at $\omega=0$). What is absent at $\omega=0$ is an energy \emph{dispersion},
not the anisotropy itself. Two consequences follow, and they set the honest scope of
the channel. First, the directional-dispersion (group-velocity) reading of this section
is a \emph{massive-sector} ($\omega\neq0$) effect, whereas the $\omega$-independent
magnitude anisotropy is a filtering/spreading effect, not an energy dispersion. Second,
a \emph{massless} energy dispersion---the object an astrophysical photon test would
constrain---is \emph{not} delivered by the $\omega\to0$ limit of the spinor propagator;
on this lattice the photon is the induced gauge field of Sec.~\ref{sec:gauge_sector},
whose
dispersion is the common-mode $\omega^2 = \mathbf{k}^{\!\top}\boldsymbol{\varepsilon}
\,\mathbf{k}$. This is the same reason the massless spinor mode is not, without a
further argument, a photon (Remark~\ref{rem:photon_limit}).

\begin{remark}[The nonlinear renormalization carries a foundational constraint]
\label{rem:nonlinear_qm}
Because the A=1 rescaling is a norm-restoring \emph{nonlinear} evolution, it inherits a
constraint that is independent of Lorentz tests. Norm-restoring nonlinear modifications of
quantum mechanics are of the postselection type and generically permit superluminal
signaling between entangled parties~\cite{weinberg1989,gisin1990,polchinski1991}:
concretely, for an entangled pair whose $A$-side branches differ in $|\mathbf{k}_\perp|$,
the universal differential damping of Sec.~\ref{subsec:as_constructed} together with a
\emph{global} renormalization can change $B$'s marginals depending on $A$'s local basis
choice, which a local completely-positive trace-preserving (CPTP) map cannot do. There is
also an internal tension to square: were the renormalization a harmless gauge, the
$\omega=0$ filtering would be unobservable, yet \texttt{exp\_01} and the analysis above
insist it \emph{is} observable. We do not resolve this here; we flag the three routes that
would, each with different phenomenology. \emph{(a)} Implement A=1 as a genuine CPTP map by
dilation, so the lost norm flows to an explicit sink sector---anisotropic \emph{absorption}
into residue rather than a nonlinear rescale, which signals nothing but opens a decay
channel with its own bounds (this dovetails the program's ``nothingness'' ontology).
\emph{(b)} A minimum-probability floor $\delta p_\text{min}$ that cuts off the differential
damping (the same regulator invoked for the filtering). \emph{(c)} An argument that the
A=1 rule, restricted to the physically preparable state space, never generates the
dangerous entangled configurations. This touches every channel that leans on the
renormalization and warrants a dedicated treatment; it is a stated open foundational item.
\end{remark}

\subsection{The as-constructed reading: an $O(1)$ matter-sector exclusion}
\label{subsec:as_constructed}

The closed forms above are exact, and it is tempting to read the resulting
anisotropy as only the small, dimension-six fractional energy shift of
Sec.~\ref{sec:falsifiability}. Taken literally, though, the bare single-tick
spinor forces three much stronger, order-unity statements about the matter sector,
all with the same geometric origin: the difference vectors
$\mathbf{V}_i-\mathbf{V}_j$ of Sec.~\ref{subsec:structure_factor} all lie in the
plane perpendicular to $\hat{\mathbf{n}}_*$, so \emph{every} axial dependence of
$\tilde H_\text{RGB}$ is a pure global phase that the modulus discards. We state
them plainly; all are verified in \texttt{exp\_05\_matter\_sector\_order}.

\begin{proposition}[Exact axial flat band]
\label{prop:flat_band}
$\theta(\mathbf{k}) = \theta(\mathbf{k}_\perp)$ to all orders, so
$\partial\theta/\partial k_\parallel \equiv 0$ at \emph{every} $\mathbf{k}$, not
merely at small $\mathbf{k}$. A wavepacket with momentum along $\hat{\mathbf{n}}_*$
does not propagate, and---since $\dot{x}_\parallel = \partial\theta/\partial
k_\parallel$ and the two-component Berry connection contributes no anomalous
velocity---$x_\parallel$ stays inert even under forces that change $k_\parallel$.
This is an exact flat band along the axis, with the attendant infinite degeneracy;
a hydrogen-like state built on this kinetic term is dynamically two-dimensional.
\end{proposition}

Two further consequences follow at fixed $|\mathbf{k}|$. First, the kinetic energy
is \emph{entirely} directional:
\begin{equation}
    \theta(\mathbf{k}) - \tfrac{\omega}{2}
      = \tfrac13\sin\omega\,|\mathbf{k}|^2\sin^2\!\alpha + O(k^4),
    \label{eq:kinetic_energy}
\end{equation}
modulated by $\sin^2\alpha$ with \emph{no} isotropic piece for the anisotropy to be
a small correction to. The dimension-six fractional number of
Sec.~\ref{sec:falsifiability} is the energy anisotropy relative to the \emph{rest}
phase $\omega/2$; the quantity a Hughes--Drever--class experiment compares---kinetic
and level structure versus orientation---is anisotropic at order unity. If the
spinor is any laboratory particle, this is excluded at zeroth order, no precision
required. Second, the scale cannot be tuned away: the lattice perpendicular kinetic
coefficient over the continuum $\mathbf{k}^2/2m$ (with $m=\sin(\omega/2)/a$) is
$\tfrac43\sin^2(\omega/2)\cos(\omega/2)\le 8/(9\sqrt3)\approx0.513$ at \emph{any}
$\omega$, so the single-tick spinor is simultaneously axially inert, $O(1)$
anisotropic, and kinetically $(ma)^2$-suppressed relative to a continuum particle.

\begin{proposition}[Linearly accumulating $\omega=0$ filtering]
\label{prop:filtering}
The eigenvalue magnitude obeys $|\lambda_+(\mathbf{k})|^2 = \sin^2(\omega/2) +
\cos^2(\omega/2)\,|\tilde H_\text{RGB}(\mathbf{k})|^2$---unity along the axis, and
$1-\tfrac43\cos^2(\omega/2)|\mathbf{k}_\perp|^2$ across it. The per-tick intensity
damping of a transverse relative to an axial component is therefore
$\tfrac43\cos^2(\omega/2)(k_\perp a)^2$; multiplied by the tick rate $c/a$, the
physical rate is $\tfrac43\cos^2(\omega/2)\,k_\perp^2\,a\,c$---\emph{linear} in $a$,
because the dimension-six per-tick suppression is repaid by the $a^{-1}$ tick count.
The dispersion channels do not accumulate (a velocity shift stays $(ka)^2$ forever);
this amplitude channel does.
\end{proposition}

At $a=\ell_P$ the resulting $e$-folding times are days to minutes---an electron at
hydrogenic momentum $\approx5$ days, a thermal $\mathrm{N}_2$ molecule
$\approx49$ minutes, a cold neutron $\approx7$ days. Even allowing for collisional
rethermalization (a gas relaxes to a steady-state momentum-distribution anisotropy
of order the filter-to-collision-rate ratio) and generous modeling slack, the effect
starts tens of orders of magnitude above co-magnetometer and matter-interferometer
bounds. The ``testability open'' status this channel formerly carried is therefore
too sanguine.

\paragraph{The honest verdict, and the escape.} Read as literal laboratory matter,
the bare single-tick spinor is \emph{excluded as constructed}: axially inert
(Prop.~\ref{prop:flat_band}), $O(1)$ anisotropic in its kinetic sector, kinematically
$(ma)^2$-mismatched to a continuum particle, and amplitude-filtering
(Prop.~\ref{prop:filtering}) at rates far above bound. This is the matter-sector
counterpart of the gauge-sector common-mode tension (Sec.~\ref{sec:gauge_sector}),
and it has the same three-of-four geometric origin. The escape is the one the program
already owns: physical matter is \emph{emergent}---composite ``session'' excitations
with their own dispersion, not the bare single-tick spinor---so the closed forms above
are the \emph{substrate} kinematics, and the observable matter dispersion (with it the
sidereal channel of Sec.~\ref{sec:falsifiability}) follows only once the
token$\to$matter map is exhibited; the program's minimum-probability floor
$\delta p_\text{min}$ is the natural regulator that halts the filtering. Until that map
is delivered, the matter-sector channels read literally are \texttt{FAIL} as
constructed, with the emergence map the open route to viability
(Table~\ref{tab:audit}). What changes is \emph{not} the closed-form
split~\eqref{eq:split}, which remains exact as substrate kinematics, but only the
claim that it registers as a small observable effect on a laboratory particle.

Equations~\eqref{eq:H2_exact}, \eqref{eq:dispersion_exact},
\eqref{eq:group_velocity} and Proposition~\ref{prop:speed_split} have been
checked against independent direct computation in
\texttt{exp\_02\_closed\_form\_split\_check}: the structure-factor and
dispersion identities to floating-point precision ($<10^{-12}$, versus the
triple sum and versus $\arg\lambda_+$ respectively), and the group
velocity and principal speed coefficients $\tfrac23\sin\omega$ and $0$ to
finite-difference truncation ($<10^{-5}$); the order-unity consequences of
Sec.~\ref{subsec:as_constructed} are verified in
\texttt{exp\_05\_matter\_sector\_order}. These checks establish the closed-form
identities (the kinematic row's \emph{Analytic} content is \texttt{PASS}); the
\emph{observable} reading of the same identities on a laboratory particle is the
as-constructed \texttt{FAIL} of Sec.~\ref{subsec:as_constructed}, pending the
emergence map. The
observability of the underlying anisotropy in renormalized amplitudes was
established separately in \texttt{exp\_01} (Table~\ref{tab:audit};
\texttt{notes/optical\_axis\_renormalization\_test.md}).

\begin{remark}[The massless mode is not automatically a photon]
\label{rem:photon_limit}
Proposition~\ref{prop:massless_flat} already shows that the $\omega\to0$ eigenphase
does not disperse. There is a second, independent reason not to read the massless
limit as a photon result: the $\omega=0$ limit of a spinor propagator is, in the
first instance, a massless \emph{spinor} mode, and identifying the two-component mode
with a photon---its helicity content, gauge character, and representation---is a
separate step this paper does not take. The robust framework-internal content of this
section is thus the closed-form uniaxial anisotropy of $|\tilde H_\text{RGB}|$ about
$\hat{\mathbf{n}}_*$ and its survival under A=1 renormalization (\texttt{exp\_01}),
with the massive-sector group velocity~\eqref{eq:split} as its clearest observable
instance; the photon-sector dispersion is treated separately in
Sec.~\ref{sec:gauge_sector}.
\end{remark}

% ── Scoping remark: mode count is unchanged by the split ──────────────
\begin{remark}[The split adds no propagating modes]
\label{rem:no_new_modes}
The axial and transverse group speeds derived above are the two
\emph{direction-dependent phases} of the bipartite propagator's existing
eigenvalue branches $\lambda_\pm(\mathbf{k})
= i\sin(\omega/2) \pm \cos(\omega/2)\,|\tilde{H}_\text{RGB}(\mathbf{k})|$.
The entire directional anisotropy enters through the single scalar
$|\tilde{H}_\text{RGB}(\mathbf{k})|$, whose low-$\mathbf{k}$ expansion
$1 - \tfrac43|\mathbf{k}_\perp|^2$ is flat along $\hat{\mathbf{n}}_*$ and
curved (coefficient $4/3$) in the perpendicular plane.  The closed-form
split is thus a modulation of the \emph{same} two branches---the
$(\psi_R,\psi_L)$ pair of the continuum Dirac spinor---and introduces no
additional branch: the mode count is two, before and after.

This is logically independent of the lattice's fermion-doubling content.
Doublers are a Brillouin-zone-corner phenomenon, whereas the
axial--transverse split is the second-order term of an expansion
about the physical mode at $\mathbf{k}=0$; the two occupy disjoint regions
of $\mathbf{k}$-space.  Paper~I deliberately defers the doubler accounting
for the bipartite $T_\text{diamond}$ Dirac operator (Paper~I~\S6.4,
\cite{menendez2026geometryfirst}) to a separate one-loop
lattice-perturbation-theory calculation,
and nothing in the present derivation depends on, alters, or settles that
count.  This $\mathbf{k}$-space disjointness argument is inherited from
Paper~I's framework and is not a new result of the present paper; in
particular, the optical-axis dispersion anisotropy does \emph{not} imply a larger
particle zoo.
\end{remark}

% ── Scoping remark: directional dispersion and packet spreading are the ──
% ── same channel read two ways, not two new falsifiers.  Folds the ────
% ── proposed post-v1 "Experiment 5" into this section per the program ─
% ── decision (4 falsifiers + 1 demonstrator; Exp 5 is the demonstrator).
\begin{remark}[Directional dispersion and wave-packet spreading restate this channel]
\label{rem:spreading_same_channel}
Two real-space signatures sometimes proposed as independent tests are, on
this lattice, the \emph{same} kinematic anisotropy read through different
observables and add no new claim. First, a direction-dependent propagation
speed---sometimes labeled ``vacuum birefringence'' in the literature---is, on
this lattice, precisely the directional dispersion anisotropy of
Proposition~\ref{prop:speed_split} and Eq.~\eqref{eq:H2_perp}, not a distinct
fixed-$\mathbf{k}$ polarization split; it is this section's subject, not a
separate prediction. Second, the anisotropic \emph{spreading} of an initially
isotropic packet is the real-space dual of the $\omega$-independent
structure-factor magnitude curvature $|\tilde H_\text{RGB}|^2 = 1 -
\tfrac43|\mathbf{k}_\perp|^2$ (Eq.~\eqref{eq:H2_perp})---\emph{not} the
massive-sector group velocity of Eq.~\eqref{eq:group_velocity}, which vanishes at
$\omega=0$---and is exactly the observable that \texttt{exp\_01} already measures
(the pulse flattens perpendicular to $\hat{\mathbf{n}}_*$, present and identical at
$\omega=0$). That this spreading is $\omega$-independent, whereas the group-velocity
contrast~\eqref{eq:split} scales as $\sin\omega$, is precisely the distinction of
Sec.~\ref{subsec:a1_normalization}: the magnitude anisotropy is a filtering/spreading
effect, the eigenphase group velocity a massive-sector energy dispersion. The isotropic-spread rate itself
($\mathrm{d}^2E/\mathrm{d}p^2$, direction-averaged) is the province of the
Lorentz-violating dispersion program, whose real-space demonstrator lives
there rather than here; only its \emph{$(1,1,-1)$-aligned anisotropy} is
this paper's result. Any claim that the lattice packet spreads ``faster than
continuum QFT'' must be stated as the residual that survives the continuum
limit $a\to0$: a free packet already spreads, and a lattice excess that
vanishes as $a\to0$ is numerical dispersion, not physics. What survives
$a\to0$ here is the \emph{direction dependence} fixed by
$\hat{\mathbf{n}}_*$, not an isotropic spread-rate enhancement.
\end{remark}
